% Couleur d'un objet opaque (absorption, diffusion) et d'un filtre (absorption, transmission) (L'essentiel, chapitre 19)
\definecolor{pcR}{RGB}{230,30,30}\definecolor{pcV}{RGB}{20,170,50}\definecolor{pcB}{RGB}{30,70,230}
\begin{tikzpicture}[x=1cm,y=1cm,
  ray/.style={line width=1.3pt,-{Stealth[length=2.2mm]}}]
% Objet opaque rouge
\begin{scope}
  \fill[pcR!70] (-0.2,-1.0) rectangle (3.0,-0.7);
  \draw[pcR!80!black,line width=0.5pt] (-0.2,-1.0) rectangle (3.0,-0.7);
  \node[pcLegende,anchor=north] at (1.4,-1.05) {objet opaque rouge};
  \draw[ray,pcR] (0.0,1.4) -- (0.75,-0.65);
  \draw[ray,pcV] (0.35,1.4) -- (1.1,-0.65);
  \draw[ray,pcB] (0.7,1.4) -- (1.45,-0.65);
  \node[pcLegende,anchor=south] at (0.35,1.45) {lumière blanche};
  \draw[ray,pcR] (1.55,-0.65) -- (2.5,1.2);
  \node[pcLegende,anchor=south,text=pcR!80!black] at (2.5,1.25) {rouge diffusé};
  \node[pcLegende,anchor=west,align=left] at (3.05,-0.85) {vert et bleu\\absorbés};
\end{scope}
% Filtre jaune
\begin{scope}[shift={(5.4,0)}]
  \fill[yellow!70!orange,opacity=0.7] (1.3,-1.0) rectangle (1.6,1.4);
  \draw[otInk!60,line width=0.5pt] (1.3,-1.0) rectangle (1.6,1.4);
  \node[pcLegende,anchor=south] at (1.45,1.42) {filtre jaune};
  \draw[ray,pcR] (-0.6,0.6) -- (1.25,0.6);
  \draw[ray,pcV] (-0.6,0.2) -- (1.25,0.2);
  \draw[ray,pcB] (-0.6,-0.2) -- (1.25,-0.2);
  \node[pcLegende,anchor=south] at (0.3,0.75) {lumière blanche};
  \draw[ray,pcR] (1.65,0.6) -- (3.3,0.6);
  \draw[ray,pcV] (1.65,0.2) -- (3.3,0.2);
  \node[pcLegende,anchor=west,align=left] at (3.35,0.4) {rouge + vert\\transmis : jaune};
  \node[pcLegende,anchor=north,text=pcB] at (1.45,-1.05) {bleu absorbé};
\end{scope}
\end{tikzpicture}
\par\smallskip\legende{La couleur perçue est celle de la lumière diffusée (objet opaque) ou transmise (filtre) : elle dépend de la lumière qui éclaire l'objet.}
